\documentclass[11pt]{article}
\usepackage[margin=1in]{geometry}
\usepackage[T1]{fontenc}
\usepackage{lmodern,amsmath,amssymb,amsthm,mathtools,hyperref}
\setlength{\emergencystretch}{3em}
\newtheorem{theorem}{Theorem}[section]
\newtheorem{lemma}[theorem]{Lemma}
\newtheorem{corollary}[theorem]{Corollary}
\newtheorem{proposition}[theorem]{Proposition}
\newcommand{\Z}{\mathbb Z}
\newcommand{\R}{\mathbb R}
\newcommand{\tr}{\operatorname{tr}}
\newcommand{\Sym}{\operatorname{Sym}}
\newcommand{\Skew}{\operatorname{Skew}}
\title{Tensor stability and arithmetic rigidity\\in Hadamard-related problems}
\author{Joshua Mark\thanks{ORCID: \href{https://orcid.org/0009-0003-5905-2382}{0009-0003-5905-2382}. AI tools assisted with research, derivation, adversarial review, computation and exposition.}}
\date{20 November 2026 --- research draft; priority not established}
\begin{document}
\maketitle
\begin{abstract}
For every finite real $p>2$, we prove that a nonflat critical point of
the entrywise $p$-energy on the real orthogonal group acquires a negative
tangent-Hessian direction at some finite self-tensor power. The proof
uses symmetric and skew-symmetric trace averages and strict Chebyshev
inequality. It gives an explicit bound on the number of tensor factors
and a factored descent direction. Consequently the orthogonal matrices
whose every tensor power is a local minimum are exactly the scaled
Hadamard matrices. This does not prove that such matrices exist in
every admissible dimension: tensor lifting changes the dimension, and
no descent back to the original dimension is established. Novelty is
not asserted pending a comprehensive literature review.
\end{abstract}

\section*{The third research pass}
The leading result of this pass is the uniform tensor-destabilization
theorem below. It extends the earlier examples to every nonflat
critical point and every finite real exponent $p>2$. The subsequent
sections prove sharp energy floors for the entire quadratic-character
template, arithmetic lower bounds on observation-trade volume, and
uniform persistence of an exact three-row augmentation barrier.
These results advance auxiliary problems. None proves all-order
Hadamard existence, and bibliographic novelty remains unestablished.


\section{Main statement}
For $n\ge2$ and $2<p<\infty$, define
\[
 F_p(U)=\sum_{i,j=1}^n |u_{ij}|^p\qquad(U\in O(n)).
\]
Call $U$ flat if all its entries have magnitude $1/\sqrt n$.
Flatness is equivalent to $\sqrt nU$ being a Hadamard matrix.
Jensen's inequality and $\sum u_{ij}^2=n$ give
\[
 F_p(U)\ge n^{2-p/2},
\]
with equality exactly when $U$ is flat. Entrywise norm optimization
and its differential framework have been studied in the almost-Hadamard
literature~\cite{BN}. Our assertion concerns stability under arbitrary
\emph{self-tensor powers}, not just the landscape at one fixed order.

\begin{theorem}[Finite tensor destabilization]\label{thm:main}
Let $U\in O(n)$ be a nonflat critical point of $F_p$. Put
\[
 q=\sum_{ij}|u_{ij}|^p,\qquad
 r=\sum_{ij}|u_{ij}|^{p-2},\qquad
 b=\frac{r+q}{(n+1)q}.
\]
Then $0<b<1$. For some positive integer $k$, the critical point
$U^{\otimes k}\in O(n^k)$ has a nonzero skew-symmetric tangent
generator $X$ such that
\[
 \left.\frac{d^2}{dt^2}F_p(U^{\otimes k}e^{tX})\right|_{t=0}<0.
\]
One may take
\begin{equation}\label{eq:kbound}
 k=2+\left\lceil\frac{\log(p-1)}{\log(1/b)}\right\rceil.
\end{equation}
The direction can be chosen in the factored form
$X=Y\otimes Z^{\otimes(k-1)}$, with $Y$ skew-symmetric and $Z$
symmetric. Both may be chosen from elementary matrix bases.
\end{theorem}

The result allows zero entries. The function $|x|^p$ is $C^2$ on the
whole real line for every $p>2$, with second derivative zero at zero.
No spectral positivity assumption on the critical-point matrix below
is needed.

\section{Derivative and trace identities}
Write
\[
 W=\bigl(|u_{ij}|^{p-2}u_{ij}\bigr)_{ij},\qquad M=U^{\mathsf T}W.
\]
Differentiation in skew directions shows that criticality is equivalent
to $M$ being symmetric. Its trace is $q$.
For any real $n\times n$ matrix $V$, set
\[
 T(V)=\sum_{ij}|u_{ij}|^{p-2}(UV)_{ij}^2,\qquad
 S(V)=\tr(M V^{\mathsf T}V).
\]
We always have $T(V)\ge0$, but individual values of $S$ may be negative.
For a skew-symmetric $V$,
\begin{equation}\label{eq:hessian}
 \left.\frac{d^2}{dt^2}F_p(Ue^{tV})\right|_0
 =p\bigl((p-1)T(V)-S(V)\bigr).
\end{equation}
Indeed the unreduced second derivative is
$p(p-1)\sum |u_{ij}|^{p-2}(UV)_{ij}^2+
p\tr(W^{\mathsf T}UV^2)$; symmetry of $M$ and
$V^2=-V^{\mathsf T}V$ give the formula.

\begin{lemma}[Symmetric and skew trace sums]\label{lem:traces}
Let $\mathcal B_+$ and $\mathcal B_-$ be Frobenius-orthonormal bases
of the symmetric and skew-symmetric matrices respectively. Then
\begin{align*}
 \sum_{Z\in\mathcal B_+}T(Z)&=(r+q)/2,&
 \sum_{Z\in\mathcal B_+}S(Z)&=(n+1)q/2,\\
 \sum_{Y\in\mathcal B_-}T(Y)&=(r-q)/2,&
 \sum_{Y\in\mathcal B_-}S(Y)&=(n-1)q/2.
\end{align*}
\end{lemma}
\begin{proof}
For the unit row vector $u_i$ of $U$, the squared norms of the
symmetric and skew projections of $u_i^{\mathsf T}e_j^{\mathsf T}$
are $(1+u_{ij}^2)/2$ and $(1-u_{ij}^2)/2$. Equivalently, direct
summation over $E_{ii}$ and $(E_{ij}\pm E_{ji})/\sqrt2$ gives
\[
 \sum_{Z\in\mathcal B_\pm}(UZ)_{ij}^2
 =\frac{1\pm u_{ij}^2}{2}.
\]
Multiply by $|u_{ij}|^{p-2}$ and sum to obtain the $T$ identities.
The same elementary bases satisfy
\[
 \sum_{Z\in\mathcal B_\pm}Z^{\mathsf T}Z
 =\frac{n\pm1}{2}I_n.
\]
Taking their trace product with $M$ proves the $S$ identities.
\end{proof}

\begin{lemma}[Strict moment inequality]\label{lem:cheb}
For a nonflat $U$, $q\le r<nq$.
\end{lemma}
\begin{proof}
The first inequality follows from $|u_{ij}|\le1$. For the second,
list the $n^2$ squared entries as $x_1,\ldots,x_{n^2}$ and put
$a=(p-2)/2>0$. Then $\sum x_i=n$, $\sum x_i^a=r$, and
$\sum x_i^{a+1}=q$. Strict Chebyshev inequality follows directly from
\[
 n^2q-nr
 =\sum_{i<j}(x_i-x_j)(x_i^a-x_j^a)>0.
\]
Every summand is nonnegative; at least one is positive because the
$x_i$ are not all equal. This argument includes zero entries.
\end{proof}

\begin{lemma}[Ratio selection without positivity of every denominator]
Suppose finitely many pairs $(T_i,S_i)$ have $T_i\ge0$ and
$\sum S_i>0$. Some $i$ has $S_i>0$ and
\[
 T_i/S_i\le\frac{\sum T_j}{\sum S_j}.
\]
\end{lemma}
\begin{proof}
Set the displayed aggregate ratio to $c\ge0$. If all positive-$S_i$
terms had $T_i>cS_i$, while nonpositive-$S_i$ terms satisfy
$T_i\ge0\ge cS_i$, summing would contradict the definition of $c$.
\end{proof}

\section{Proof of tensor destabilization}
Apply the trace sums and ratio-selection lemma to choose $Y\in
\mathcal B_-$ and $Z\in\mathcal B_+$ with positive $S$-values and
\[
 \alpha:=\frac{T(Y)}{S(Y)}
 \le\frac{r-q}{(n-1)q}<1,\qquad
 \beta:=\frac{T(Z)}{S(Z)}\le b<1.
\]
The upper bound on $\alpha$ is nonnegative, and the choices exist
because $n\ge2$ and $q>0$.

At the tensor power, $W_k=W^{\otimes k}$ and $M_k=M^{\otimes k}$.
Hence $U^{\otimes k}$ is still critical. For
$X=Y\otimes Z^{\otimes(k-1)}$, which is nonzero and skew-symmetric,
tensor factorization of the sums and traces gives
\[
 T_k(X)=T(Y)T(Z)^{k-1},\qquad
 S_k(X)=S(Y)S(Z)^{k-1}.
\]
Its Hessian is therefore
\begin{equation}\label{eq:tensorhessian}
 pS(Y)S(Z)^{k-1}
 \bigl((p-1)\alpha\beta^{k-1}-1\bigr).
\end{equation}
If $\alpha=0$, $k=1$ already works. Otherwise, choose $k$ as in
\eqref{eq:kbound}. Since $\alpha<1$, $\beta\le b<1$, and
$k-1>\log(p-1)/\log(1/b)$, the bracket is negative. This proves
Theorem~\ref{thm:main}.

\begin{corollary}[Characterization by tensor stability]
For $p>2$, the following are equivalent for $U\in O(n)$:
\begin{enumerate}
\item Every positive self-tensor power of $U$ is a local minimum of
the entrywise $p$-energy in its ambient orthogonal group.
\item Every positive self-tensor power is critical with positive
semidefinite tangent Hessian.
\item $U$ is flat, so $\sqrt nU$ is a Hadamard matrix.
\end{enumerate}
\end{corollary}
\begin{proof}
The first implication is the second-order necessary condition. The
theorem proves the second implication to flatness. If $U$ is flat,
every tensor power is flat and attains the Jensen lower bound globally.
\end{proof}

\begin{corollary}[Strict eventual submultiplicativity]
Let $m_p(n)=\min_{U\in O(n)}F_p(U)$. If there is no Hadamard matrix
of order $n$, then for some positive $k$,
\[
 m_p(n^k)<m_p(n)^k.
\]
Equivalently, equality for every $k$ occurs precisely when order $n$
admits a Hadamard matrix.
\end{corollary}
\begin{proof}
Compactness gives a minimizer $U$, which is critical and nonflat under
the hypothesis. The negative Hessian direction strictly decreases
$F_p$ near $U^{\otimes k}$, whose value is $m_p(n)^k$. Conversely a
flat matrix and all its tensor powers attain the Jensen lower bounds.
\end{proof}

\begin{corollary}[Rectangular row-orthonormal frames]
The tensor-destabilization theorem and tensor-stability characterization
also hold on
\[
 \operatorname{St}(d,n)=\{U\in\R^{d\times n}:UU^{\mathsf T}=I_d\},
 \qquad 1\le d\le n,\quad n\ge2.
\]
Here flat means every entry has magnitude $1/\sqrt n$, so the scaled
matrix is a rectangular partial Hadamard matrix. The tensor powers have
shape $d^k\times n^k$, and the same bound~\eqref{eq:kbound} uses the
number $n$ of columns.
\end{corollary}
\begin{proof}
Right multiplication by $O(n)$ acts transitively on this Stiefel
manifold, so criticality is again equivalent to symmetry of the
$n\times n$ matrix $M=U^{\mathsf T}W$. All rows still have norm one,
so both trace-sum identities are unchanged. For strict Chebyshev, there
are now $dn$ squared entries summing to $d$; the same pair sum is
$dnq-dr>0$ unless every entry is flat. Thus $r<nq$ still holds and
the ratio and tensor arguments apply without change. Positive
$S(Y),S(Z)$ ensure that $UY$ and $UZ$ are nonzero. Indeed, if
$UV=0$ for a symmetric or skew $V$, then $VU^{\mathsf T}=0$ and
$V$ annihilates the range of $M$, implying $S(V)=0$.
The factored actual tangent is therefore nonzero. Finally flat row frames minimize the energy
globally, by Jensen separately on each unit row, with value
$d n^{1-p/2}$.
\end{proof}

\section{A local minimum that becomes unstable after two copies}
For $p=4$, let $A=2J_5/5-I_5$. The earlier note proves that
$A\otimes(H_k/\sqrt k)$ is a strict local minimum with value $29/25$
for any Hadamard factor. At $A$,
\[
 M=\frac7{25}I_5-\frac6{125}J_5.
\]
Let $Y=E_{01}+E_{12}-E_{02}$ be a skew triangle circulation. Let $Z$
be symmetric, with $Z_{01}=1$, $Z_{12}=-1$, $Z_{23}=1$,
$Z_{30}=-1$, the transposed entries equal, and all other entries zero.
Both have row sums zero and zero diagonal, and
$\|Y\|_F^2=6$, $\|Z\|_F^2=8$. Thus
\[
 T(Y)=\tfrac4{25}\|Y\|_F^2,\quad
 S(Y)=\tfrac7{25}\|Y\|_F^2,
\]
with the same identities for $Z$. Consequently at $A\otimes A$ the
direction $Y\otimes Z$ has quartic curvature
\[
 4\cdot\frac{49}{625}\cdot48
 \left(3\frac{16}{49}-1\right)=-\frac{192}{625}.
\]
Tensoring with a flat factor preserves this displayed curvature when
the tangent generator is tensored with an identity. In particular a
strict local minimum at order $20$ becomes unstable after taking its
tensor square at order $400$. This is consistent with stability under
flat tensor factors: self-tensoring inserts a second nonflat factor.

\section{Scope, checks and priority}
The proof is uniform in all real finite $p>2$, not a claim inferred
from numerical Hessian spectra. The accompanying standard-library
Python checker verifies the trace formulas, stationarity, ratio choices,
and factored negative curvatures with exact rational arithmetic for
integer exponents $p=3,4,5,6,8$ and several critical matrices. It also
checks the two-copy curvature directly at order $25$.

A separate mathematical review checked the trace formulas, strictness
with zeros, possibly indefinite $M$, and the tensor factorization.
This is not a Lean proof or a substitute for external peer review.
Targeted primary-literature searches on almost-Hadamard matrices,
tensor stability, local norm extrema and tensor-power descent did not
locate this precise statement. That does not establish priority.

Most importantly, this result does not show that a tensor-stable local
minimum exists in any prescribed dimension. Tensor descent occurs in
$O(n^k)$, and no mechanism here transfers it back to $O(n)$. Therefore
the characterization is an optimization result and a research lead,
not a proof of the Hadamard conjecture.


% Reviewed branch fragment; requires theorem/proposition/proof environments.
\section{The exact energy floor of the full quadratic-character template}
The earlier constructions used one particular arrangement of a quadratic
character, constants, and the all-one polynomial. The following result
shows that their minimum energies cannot be improved by taking arbitrary
integer linear combinations of those same three ingredients, even with
independent shifts in the four channels. This is a restriction on a
construction template, not on unrestricted Hadamard existence. No claim
of bibliographic novelty is made.

Let $m>1$ be odd, $R_m=\Z[z]/(z^m-1)$, $z^*=z^{-1}$ and
$J=1+z+\cdots+z^{m-1}$. Suppose $L$ has coefficient zero at the
identity, all other coefficients in $\{\pm1\}$, and
\[
 LL^*=m\cdot1-J,\qquad L^*=\eta L,\qquad \eta\in\{+1,-1\}.
\]
For prime $m=p$, the quadratic character polynomial has these
properties with $\eta=(-1)^{(p-1)/2}$. Its coefficients take both signs:
augmentation of the norm identity forces $L(1)=0$.
Consider the four-channel family
\begin{equation}\label{eq:fullspan}
 A_i=z^{r_i}(u_iL+v_i\cdot1+w_iJ),\qquad
 u_i,v_i,w_i,r_i\in\Z.
\end{equation}
Require every coefficient of every $A_i$ to be odd and
$\sum_iA_iA_i^*=4t\cdot1$.

\begin{theorem}[Sharp floor in the full character span]\label{thm:spanfloor}
If $m\equiv3\pmod4$ and $L^*=-L$, the least possible energy in
\eqref{eq:fullspan} is
\[
 T_-(m)=\frac{m^2+6m+21}{16}.
\]
If $m\equiv1\pmod4$, $m\ge5$, and $L^*=L$, the least possible energy is
\[
 T_+(m)=\frac{m^2+2m+5}{8}.
\]
Both minima are attained by explicit odd-coefficient choices. The
integers $T_\pm(m)$ are odd. These are global minima within the whole
family, not only within the earlier special constructions.
\end{theorem}
\begin{proof}
Shifts do not change the individual Hermitian norms and preserve
coefficient parity, so omit them. Since the nonidentity coefficients
of $L$ include both signs, coefficientwise oddness is equivalent to
\begin{equation}\label{eq:spanparity}
 u_i\equiv v_i\equiv1-w_i\pmod2.
\end{equation}
Put $S=\sum_i u_i^2$, $V=\sum_i v_i^2$, and $W=\sum_i w_i^2$.
Since $LJ=0$, the norm sum is
\[
 (mS+V)\cdot1+
 \left(-S+2\sum_iw_iv_i+mW\right)J
 +(1+\eta)\left(\sum_i u_iv_i\right)L.
\]
The three polynomials $1,J,L$ are linearly independent over $\mathbb Q$.
Consequently scalarity requires
\begin{equation}\label{eq:spanconditions}
 2w\mathbin{\cdot}v=S-mW,\qquad 4t=mS+V,
\end{equation}
and in the symmetric case also $u\mathbin{\cdot}v=0$.
The energy $t$ is odd, since $4m$ odd coefficients have squared sum
$4$ modulo eight. Also $W>0$: otherwise scalarity forces $S=0$,
contradicting~\eqref{eq:spanparity}. Cauchy--Schwarz gives the useful
universal inequality
\begin{equation}\label{eq:spanCS}
 t\ge\frac{mS+(S-mW)^2/(4W)}4
   =\frac{(S+mW)^2}{16W}.
\end{equation}

In the skew case, first suppose $W\ge2$. Then $t\ge m^2/8$, while
\[
 \frac{m^2}{8}-(T_-(m)-2)
 =\frac{(m-3)^2+2}{16}>0.
\]
Since $t$ and $T_-(m)$ are odd integers, this implies $t\ge T_-(m)$.
If $W=1$, exactly one $w_i$ is $+1$ or $-1$ and the other three
vanish. Thus three $u_i$ and three corresponding $v_i$ are odd,
giving $S\ge3$ and
\[
 V\ge\frac{(S-m)^2}{4}+3.
\]
It follows that
\[
 t\ge\frac{mS+(S-m)^2/4+3}{4}
 =\frac{(m+S)^2}{16}+\frac34
 \ge T_-(m).
\]
For attainment write $m=4k+3$, and take
\[
 (A_1,A_2,A_3,A_4)=(L+1,L+1,L+1,J-2k).
\]
The first three norm sums contribute $(3m+3)1-3J$, while the last
contributes $4k^2\cdot1+3J$, giving energy $T_-(m)=k^2+3k+3$.

In the symmetric case, $u\cdot v=0$ and
\eqref{eq:spanparity} imply that the number of odd $w_i$ is even.
Hence $W$ is a positive even integer. If $W\ge4$,
\eqref{eq:spanCS} gives $t\ge m^2/4$. Since
\[
 \frac{m^2}{4}-(T_+(m)-2)
 =\frac{(m-1)^2+10}{8}>0,
\]
odd integrality again gives $t\ge T_+(m)$. If $W=2$, exactly two
$w_i$ are $\pm1$; the other two vanish. Their corresponding $u_i,v_i$
are odd, so $S\ge2$. Applying Cauchy--Schwarz only on the two
nonzero-$w_i$ coordinates yields
\[
 V\ge\frac{(S-2m)^2}{8}+2.
\]
Thus
\[
 t\ge\frac{mS+(S-2m)^2/8+2}{4}
 =\frac{(S+2m)^2}{32}+\frac12\ge T_+(m).
\]
For attainment, write $m=4k+1$ and set
\[
 (A_1,A_2,A_3,A_4)=(L+1,L-1,J-2k,J-2k).
\]
The two $L$ norms have cancelling $L$ terms; their $J$ coefficient
is $-2$, while the last two norms have $J$ coefficient $+2$.
The scalar energy is $2k^2+2k+1=T_+(m)$.
\end{proof}

\begin{corollary}[Only three prime lengths can reach the binary boundary]
For the quadratic character polynomial at an odd prime $p$, the full
family~\eqref{eq:fullspan} reaches energy $t=p$ if and only if
$p\in\{3,5,7\}$. In particular no prime $p\ge11$ can be settled
by rescaling, independently shifting, or taking arbitrary integer
linear combinations of a quadratic character, its zero-position
monomial, and $J$ in the four channels.
\end{corollary}
\begin{proof}
For skew lengths,
$T_-(p)-p=(p-3)(p-7)/16$; for symmetric lengths,
$T_+(p)-p=(p-1)(p-5)/8$. The displayed attaining constructions
reach the boundary at $3,5,7$ and have larger minimum energy at all
other odd primes. At the boundary every coefficient is a sign.
\end{proof}

This also covers affine relabelings of the quadratic character:
$\chi(aj+b)$ is, up to sign, a shift of $\chi(j)$, with its unique
zero shifted at the same time. It does not cover arbitrary changes
at additional positions, nonquadratic characters, or general sequence
quadruples. In particular it is not an obstruction to Hadamard matrices
of the corresponding orders.

\subsection{Fixing only two character channels already excludes infinitely many primes}
Let $p\equiv3\pmod4$, and let two channels be the binary sequences
$z^r(L+\varepsilon)$ and $z^s(L+\delta)$, where
$\varepsilon,\delta\in\{\pm1\}$. Their row sums have squares one.
If two otherwise unrestricted binary channels complete these to a
four-sequence complementary set, augmentation forces their row sums
$a,b$ to satisfy
\[
 a^2+b^2=4p-2.
\]
Both are odd; equivalently
$((a+b)/2)^2+((a-b)/2)^2=2p-1$.
For $p\equiv11\pmod{36}$, the integer $2p-1$ is three modulo nine,
so its $3$-adic valuation is exactly one. A sum of two squares cannot
have that valuation: divisibility by three forces both summands to
be divisible by three, making the valuation at least two. Thus no
such completion exists at these primes, regardless of how the other
two channels are chosen. Dirichlet's theorem supplies infinitely many
primes in the class $11\pmod{36}$; an English translation of the
original memoir is available at \url{https://arxiv.org/abs/0808.1408}.

Fixing three binary character channels imposes the stronger necessary
condition that $4p-3$ be a square. With only one fixed channel, the
remaining augmentation equation is a sum of three odd squares equal
to $4p-1\equiv3\pmod8$ and has no obstruction from the three-square
theorem. This does not assert that the full sequence constraints are
then solvable.

\paragraph{Scope and reproduction.}
The exact checker constructs the attaining matrices for primes through
$101$, checks their periodic autocorrelations, samples scalar templates
with arbitrary integer coefficients, and verifies the two-channel
congruence obstruction at small primes. The infinite assertions rest
on the proofs, not on these tests. These restrictions identify the
precise need to leave the quadratic-character affine span; they do not
improve the main Hadamard exponent bound or prove new all-order existence.

% Research fragment: theorem/proposition/corollary/proof, amsmath/amssymb/hyperref.
\section{Arithmetic rigidity of saturated Hadamard trades}
This section proves an obstruction to bounded-degree positive moves on the
\emph{balanced saturated fiber}, rather than on arbitrary binary marginal
fibers. General unbounded Markov width for binary complete-graph models is
already known, with much stronger global bounds; see Develin--Sullivant,
\url{https://arxiv.org/abs/math/0308280}, and
Kr\'al'--Norine--Pangr\'ac,
\url{https://arxiv.org/abs/0810.1979}. No bibliographic novelty is claimed
for the specialized deductions below.

Normalize a Hadamard matrix $H$ of order $n$ so its first column is all
ones. Deleting that column and replacing signs by bits gives a binary
strength-two array with $n$ observations and $d=n-1$ coordinates.
A \emph{reduced observation trade of volume $r$} deletes $r$ old
observations and inserts $r$ new ones, preserves all moments of degree at
most two, and has disjoint deleted and inserted pattern multisets.
Every array in this fiber is again a normalized Hadamard matrix: its
$n$ columns are orthogonal, so its $n$ rows are orthogonal too.
Consequently a trade preserves $n-r$ common orthogonal rows.

\begin{theorem}[Denominator obstruction to a new sign row]\label{trd:denom}
Let $n=4m$ with $m>1$ odd, and let $p$ be the least prime factor of $m$.
Suppose a sign vector $x$ lies in the real span of $r<n$ rows
$h_1,\ldots,h_r$ of an order-$n$ Hadamard matrix, but is not any
$\pm h_i$. Then
\begin{equation}\label{trd:bound}
 p\le\frac{r^{r/2}}{2^{r-1}}.
\end{equation}
More precisely, the least common denominator of the expansion
coefficients of $x$ divides $m$ and divides the normalized determinant
$\det S/2^{r-1}$ of every invertible $r\times r$ sign-column anchor $S$.
\end{theorem}
\begin{proof}
Write $x=\sum_i a_i h_i$. Orthogonality gives
$a_i=\langle x,h_i\rangle/n$ and $\sum_i a_i^2=1$.
Choose a Hadamard row outside this span and multiply columns so that row
becomes all ones. Then $x$ and all $h_i$ are balanced. Inner products of
balanced sign vectors of length divisible by four are multiples of four:
if their positive supports meet in $b$ positions, the inner product is
$4b-n$. Hence $a_i\in m^{-1}\mathbb Z$.

Let $q$ be the least positive common denominator of the $a_i$.
If $q=1$, the unit-norm integer vector $(a_i)$ is a signed coordinate
vector, contrary to the hypothesis. Thus $q>1$ and $q\mid m$.
The $r$ rows are independent, so choose $r$ column positions forming an
invertible sign matrix $S$. The equations at those positions and Cramer's
rule give $q\mid\det S$. Every sign determinant of order $r$ is divisible
by $2^{r-1}$, by subtracting one row from all the others. Since $q$ is
odd, $q\mid\det S/2^{r-1}$. Hadamard's determinant inequality now gives
\[
 p\le q\le\frac{|\det S|}{2^{r-1}}
 \le\frac{r^{r/2}}{2^{r-1}}.
\]
The same reasoning applies to every invertible anchor.
\end{proof}

\begin{corollary}[A prime-factor lower bound for positive trade volume]
Every nonempty reduced observation trade of volume $r<n$ from such an
array satisfies \eqref{trd:bound}. In particular, if
\[
 p>\max_{1\le k\le R}\frac{k^{k/2}}{2^{k-1}},
\]
then every array in this saturated fiber is isolated under all positive
moment-preserving moves of degree at most $R$.
\end{corollary}
\begin{proof}
Every inserted row is perpendicular to the $n-r$ retained rows, so lies
in the span of the $r$ deleted rows. A reduced nonempty trade contains
a genuinely new row. Its first coordinate is one, so it cannot be the
negative of an old normalized row. Apply Theorem~\ref{trd:denom}.
\end{proof}
For example, Paley order $212$ has $m=53$. Since
$8^8<53^2\,4^7$, no reduced trade of volume at most eight can leave
\emph{any} array in its saturated fiber. The theorem does not require
that the chosen array have Paley symmetry; Paley only supplies existence.

\subsection{Infinitely many existing fibers require unbounded move degree}
\begin{theorem}[Uniform bounded-degree moves cannot suffice]\label{trd:unbounded}
For every fixed $R$, there are infinitely many known Hadamard orders
$n\equiv4\pmod8$ for which every saturated array is isolated under all
nonzero moment-preserving moves of degree at most $R$.
\end{theorem}
\begin{proof}
Choose an integer $B$ larger than the determinant bound in the preceding
corollary, and put
\[
 P_B=\prod_{3\le\ell\le B\atop\ell\text{ prime}}\ell,
 \qquad M_B=8P_B.
\]
The Chinese remainder theorem gives a reduced residue class satisfying
$q\equiv3\pmod8$ and $q\equiv1\pmod\ell$ for every odd prime
$\ell\le B$. Dirichlet's theorem gives infinitely many primes $q$ in
this class. Paley's construction gives a Hadamard matrix of order
$n=q+1$. Here $m=(q+1)/4$ is odd and has no prime factor at most $B$,
because $q+1\equiv2\pmod\ell$. For sufficiently large $q$, $m>1$.
The determinant obstruction applies at every degree at most $R$.
\end{proof}

The theorem concerns these particular minimum-run, balanced fibers. It
is not presented as a new proof of the much stronger known global
Markov-width lower bounds. It also does not refute bounded row-exchange
\emph{augmentation} for arbitrary partial matrices: a trade between two
complete arrays and an augmentation of an incomplete array are different
operations. Global symmetry operations can bypass this degree restriction
by changing many observations at once.

\begin{corollary}[A quantitative infinite subsequence]\label{trd:linnik}
There is an infinite sequence of known Hadamard orders $n\to\infty$
such that every nonempty reduced observation trade has volume at least
\[
 (2-o(1))\frac{\log\log n}{\log\log\log n}.
\]
Equivalently, for every $\varepsilon>0$, this lower bound holds with
$2-\varepsilon$ along all sufficiently large orders in that sequence.
\end{corollary}
\begin{proof}
For each $B$, choose the least prime $q_B$ in the CRT class above.
Linnik's theorem gives $q_B\le C M_B^L$ for absolute constants $C,L$.
Only the existence of such constants is used; no numerical value of
Linnik's exponent is needed. The elementary Chebyshev bound gives
$\log M_B=O(B)$, hence $\log(q_B+1)=O(B)$.
Also $q_B>B$, so the selected primes tend to infinity as $B$ does.
Set $n_B=q_B+1$. The least prime factor of $(q_B+1)/4$ exceeds
$B\ge c\log n_B$ for an absolute $c>0$.

For $r\le(2-\varepsilon)\log\log n_B/\log\log\log n_B$,
\[
 \log\!\left(\frac{r^{r/2}}{2^{r-1}}\right)
 =\frac r2\log r-(r-1)\log2
 \le(1-\varepsilon/2+o(1))\log\log n_B
 <\log B.
\]
The denominator obstruction excludes every such trade. Passing to
distinct orders if necessary gives the asserted sequence.
\end{proof}
Linnik's least-prime bound is stated in the introduction of Xylouris,
\emph{On the least prime in an arithmetic progression and estimates for
the zeros of Dirichlet $L$-functions}, Acta Arith. 150 (2011), 65--91,
\url{https://doi.org/10.4064/aa150-1-4}. For the elementary bound
$\log M_B=O(B)$, one may bound products of primes in $(u,2u]$ by
$\binom{2u}{u}<4^u$ and sum over dyadic intervals.

\subsection{These fibers have many states, and a universal large trade}
\begin{proposition}[A universal half-row trade]
Every normalized Hadamard matrix of order $n>4$ has a reduced
observation trade of volume $n/2$.
\end{proposition}
\begin{proof}
Choose a nonconstant column and let $I$ be its $k=n/2$ positive rows.
On these rows the distinguished all-one column and the chosen column
are constant. Each other column has sum zero on $I$, by its
orthogonality to both columns. Apply $Q=2J_k/k-I_k$ to the selected
rows. This preserves the two constant columns and negates every other
selected column, so the entries stay signs and the first column stays
one. The transformation is orthogonal and leaves all other rows fixed,
so the resulting matrix is Hadamard. Since $k>2$, each row of $Q$
has more than one nonzero coordinate. Thus no new row is an old row
up to sign; after cancellation the trade really has volume $k$.
\end{proof}

\begin{proposition}[Exponential multiplicity of isolated states]
If $n>4$ and $n\equiv4\pmod8$, coordinatewise bit translations of
any saturated array produce $2^{n-1}$ distinct arrays in the same fiber.
For the orders in Theorem~\ref{trd:unbounded}, all are isolated under
the indicated bounded-degree move family.
\end{proposition}
\begin{proof}
First, a closed quadruple of Hadamard rows is impossible at these
orders. If four rows have pointwise product constant, column-sign
normalization and orthogonality partition columns into four fields
of size $n/4$. Orthogonality of any additional row to all four forces
its sum in each field to be zero. Thus $n/4$ must be even, contrary
to $n\equiv4\pmod8$.

Coordinatewise bit translation preserves the balanced degree-two
moments. If a nonzero translation $a$ stabilized the array, choose
two distinct translation pairs $\{x,x+a\}$ and $\{y,y+a\}$. Their four
sign rows have pointwise product one, making a closed quadruple.
Since there are $n>4$ rows, such pairs exist. Hence the translation
stabilizer is trivial, and all $2^{n-1}$ translations are distinct.
The isolation statement follows from the previous theorem.
\end{proof}

These deductions concern local moves on existing objects. They do not
prove existence at an unknown order. They show that a successful
positive-conversion or sampling argument cannot rely on a degree bound
independent of the dimension, even on the highly symmetric balanced
saturated targets. The exact scripts check the low-dimensional anchor
classification, explicit Paley examples, and small positive trades; the
infinite conclusions rest on the proofs above.

\section{Tensor products preserve a three-row exchange barrier}
For a partial Hadamard matrix $P$, let $\rho(P)$ be the minimum number
of old rows deleted in an exchange which inserts one more row than it
deletes, all retained and inserted rows remaining mutually orthogonal.
Set $\rho(P)=\infty$ if there is no such exchange. The previous pass
certified $\rho(P)=3$ for one $16\times24$ matrix. Here we prove a
uniform preservation theorem, rather than search larger isolated orders.
No claim of bibliographic novelty is made.

\begin{proposition}[Tensoring cannot amplify the exchange cost upward]
If $H$ is a full Hadamard matrix, then
\[
 \rho(P\otimes H)\le\rho(P).
\]
\end{proposition}
\begin{proof}
Choose any row $h$ of $H$. A $k$-for-$(k+1)$ exchange of rows $p_i$
for rows $v_j$ lifts to an exchange of $p_i\otimes h$ for $v_j\otimes h$.
Rows in other Hadamard fibers are orthogonal to every inserted row,
while the same-fiber checks reduce to the original exchange.
\end{proof}
Thus repeated full-Hadamard tensoring alone cannot turn a finite
exchange lower bound into an unbounded one. The interesting issue is
whether the tensor permits a cheaper exchange by using non-product rows.

\subsection{A condition excluding cheaper non-product exchanges}
For a row-index set $I$, let $\mathcal C_I$ consist of the sign vectors
orthogonal to all rows of $P$ outside $I$. Sign vectors differing only
by a global minus sign are treated as one line. Write the order as $n$.

\begin{proposition}[A general span-rigidity criterion]
Suppose $P$ is maximal and, for every index set $I$ of size at most $c$,
\[
 \mathcal C_I\subseteq\operatorname{span}\{p_i:i\in I\}.
\]
Then no tensor $P\otimes H$ with a full Hadamard factor admits an
augmenting exchange deleting at most $c$ rows.
\end{proposition}
\begin{proof}
For a deletion set $D$ of size $k\le c$, let $I$ be its projection
onto the original row indices. Reshape a candidate as $X$. As in the
proof below, every column belongs to $\mathcal C_I$, so $X$ lies in
$\operatorname{span}(p_i:i\in I)\otimes\R^s$. Orthogonality to the
retained tensor rows then forces $X$ into the span of the $k$ deleted
tensor rows. That space has dimension $k$ and cannot contain $k+1$
mutually orthogonal nonzero rows.
\end{proof}
The order-twenty-four example fails this criterion at some pairs: its
two extra sign lines span a third direction. Their particular geometry,
rather than span rigidity, supplies the stronger preservation below.

\begin{theorem}[Preservation under half-overlap pair structure]
Suppose $P$ is maximal and has the following properties:
\begin{enumerate}
\item $\mathcal C_{\{i\}}$ has just the line of $p_i$.
\item For every pair $i,j$, either $\mathcal C_{\{i,j\}}$ has only the
lines of $p_i,p_j$, or it has exactly four lines. In the latter case
the two extra lines can be oriented as $x,y$ so that
\[
 x\cdot y=0,\qquad
 (p_i\cdot x,p_j\cdot x)=(p_i\cdot y,p_j\cdot y)
 =\tfrac n2(1,\sigma),\quad\sigma\in\{\pm1\}.
\]
\end{enumerate}
Then for every full Hadamard matrix $H$, $P\otimes H$ has no augmenting
exchange deleting at most two rows. In particular, if $\rho(P)=3$, then
$\rho(P\otimes H)=3$.
\end{theorem}
\begin{proof}
Let $H$ have order $s$, with rows $h_b$, and reshape a proposed sign
row of length $ns$ as an $n\times s$ sign matrix $X$. Its correlations
with the old tensor rows are the entries of $PXH^{\mathsf T}$.
If the deleted tensor rows project to original row-index set $I$,
every column of $X$ belongs to $\mathcal C_I$: for $k\notin I$,
$p_kXH^{\mathsf T}=0$ implies $p_kX=0$.

With no deletions, this contradicts maximality. With one deletion
$(i,b)$, every column is $\pm p_i$, so $X=p_i\otimes\varepsilon$;
the remaining fiber constraints force $\varepsilon$ to be proportional
to $h_b$. The only candidate line is the removed tensor row.

If both deleted rows have the same original index $i$ and different
fibers $b,c$, the same argument gives $X=p_i\otimes\varepsilon$, now
with $\varepsilon\in\operatorname{span}(h_b,h_c)\cap\{\pm1\}^s$.
That intersection has only $\pm h_b,\pm h_c$. Indeed write
$\varepsilon=\alpha h_b+\beta h_c$; the two orthogonal sign rows
agree and disagree on nonempty sets, so both $\alpha+\beta$ and
$\alpha-\beta$ are signs. Consequently one coefficient is zero and
the other is a sign. There are at most two candidate lines.

It remains to delete $(i,b),(j,c)$ with $i\ne j$. Fourier inversion
along the Hadamard factor gives
\[
 p_iX=\alpha h_b,\qquad p_jX=\beta h_c.
\]
For each column of $X$, the possible correlation pairs with $p_i,p_j$
are $\pm(n,0)$, $\pm(0,n)$ and, when present,
$\pm(n/2)(1,\sigma)$. Their absolute magnitudes are constant across
columns, because $h_b,h_c$ have sign entries. If one of $\alpha,\beta$
vanishes, the candidate is one of the removed old tensor rows. If both
are nonzero, the extra pair type is necessary, and its fixed correlation
ratio forces $h_b$ and $h_c$ to be proportional. Thus $b=c$; otherwise
there are again only the two removed candidate lines.

In the remaining case put $h=h_b=h_c$ and orient globally so that the
extra correlation pair is $(n/2)(1,\sigma)h$.
Let
\[
 w=(p_i+\sigma p_j)/2,\qquad z=(x-y)/2.
\]
Projection onto $\operatorname{span}(p_i,p_j)$ shows that
$x=w+z$, $y=w-z$, $w\perp z$, and
$\|w\|^2=\|z\|^2=n/2$. To justify the residual cancellation, the
residuals of $x,y$ each have squared norm $n/2$ and inner product
$-n/2$, so they are negatives of each other. Every extra tensor
candidate therefore has the form
\[
 W_{\varepsilon}=w\otimes h+z\otimes\varepsilon,
 \qquad\varepsilon\in\{\pm1\}^s.
\]
Neither removed old tensor row is orthogonal to any $W_\varepsilon$.
For two extra candidates,
\[
 W_\varepsilon\cdot W_\delta
 =\tfrac n2(s+\varepsilon\cdot\delta),
\]
which is zero exactly when $\delta=-\varepsilon$. Thus the candidate
orthogonality graph consists of the old orthogonal pair and a matching
on the extra candidates. It has no triangle, so no three mutually
orthogonal replacement rows exist. All cases are exhausted.
The upper bound from the preceding proposition proves equality when
$\rho(P)=3$.
\end{proof}

\subsection{An infinite family with exact exchange cost three}
The independently certified $16\times24$ matrix from the previous pass
satisfies the hypotheses. Among its $120$ two-row deletions, one hundred
have only the two old candidate lines. The other twenty have exactly
two extra lines. Those extras are mutually orthogonal, their
correlations with the two old rows are all $\pm12$, and the two
correlation vectors are parallel. All these identities are checked
exactly from the exhaustive sign-cube certificate.

\begin{corollary}
For every Hadamard order $s$, there is a maximal partial matrix of
shape $16s\times24s$ whose minimum augmenting exchange deletes exactly
three rows. In particular such examples exist at every order
$24\cdot2^a$.
\end{corollary}

This is a statement about partial matrices and the cost of escaping
their terminal configurations. It does not produce new Hadamard orders.
The base three-for-four exchange has an explicit full completion. Apply
that exchange separately in each of the $s$ Hadamard fibers. After
$s$ exchanges the result is $Q\otimes H$, where the base $17$-row
matrix $Q$ completes; hence this whole tensor family can be completed
after those $s$ exchanges. A single lifted exchange is only claimed
to augment the tensor partial matrix, not necessarily complete it.

\paragraph{Exact checks.}
The script \path{round3/exchanges/check.py} re-enumerates all $2^{23}$
normalized sign vectors for the base matrix, verifies the pair-structure
hypotheses for all $120$ pairs, and checks lifted exchanges and
completions at tensor-factor orders $1,2,4,8$. Additional exact
low-dimensional nullspace enumeration at factor order two confirms
the predicted candidate sets in each of the three deletion cases.
The uniform statement rests on the proof above. No Lean verification
or novelty claim is made.

\section{What these results change}
The tensor theorem provides an exact structural characterization:
only scaled Hadamard matrices remain local minima at every self-tensor
power. It does not establish a tensor-stable point in any given
dimension or a way to bring a lifted improvement back to that dimension.
That is the missing step in any proposed Hadamard-existence application.

The character-span theorem closes a whole tuning direction: integer
rescaling and independent shifts cannot lower its energy floor.
Additional sequence structure, rather than a different choice of the
same constants, is necessary beyond the small exceptional lengths.

The trade bounds concern reduced observation trades on complete
balanced arrays. They do not exclude global symmetry operations,
unbounded-degree moves, or bounded augmentations of incomplete arrays.
Conversely, tensor persistence of exchange cost three is an exact
infinite family, not a proof that exchange costs are unbounded.

All source and code for this pass is under
\path{research/full_conjecture/round3/}. The four branch checkers use
only Python's standard library. The mathematical proofs, rather than
finite computations, establish the uniform assertions. No Lean code,
Git commits, submissions or publications were made in this pass.
\begin{thebibliography}{1}
\bibitem{BN} T. Banica and I. Nechita,
\emph{Almost Hadamard matrices: the case of arbitrary exponents},
Discrete Appl. Math. 161 (2013), 2367--2379.
\url{https://arxiv.org/abs/1211.2669}.
\end{thebibliography}
\end{document}
